\SourceLeaf{26}
10. If $\xi$ is a prime polynomial mod. $(1+x+xx)$ and its determinant is $=D$,
then in general

\[
\zeta^D \equiv \zeta \quad(\mathrm{mod.}\ 1+x+xx,\ \xi).
\]

11. Quadratic and cubic residues modulo $\ldots$

12. If $P$ is a cubic residue of $\xi$, then:

\[
P^{\frac{1}{3}(D-1)} \equiv 1 \quad(\mathrm{mod.}\ 1+x+xx,\ \xi).
\]



\Needspace{8\baselineskip}\begingroup\small
\begin{center}
\textbf{REMARKS ON THE DEVELOPMENTS »ON THE THEORY OF CUBIC} 
\textbf{RESIDUES«.}
\end{center}
\noindent\hspace*{\parindent}\textsc{Gauss} attempted here to extend his third proof of
the law of quadratic reciprocity to cubic residues. The first approach concerns
only the domain of rational integers and uses the following basic idea:

If the prime number satisfies $p \equiv 1 \quad(\mathrm{mod.}\ 3)$, then the congruence $x^3 \equiv 1
\quad(\mathrm{mod.}\ p)$
has three distinct roots. Let one that differs from $1$ be called $f$.
In the system of all $p-1$ integers incongruent modulo $p$ and relatively prime to
$p$, the three numbers $1, f, f^2$ form, under multiplication, a subgroup of index
$\dfrac{p-1}{3}$. For this subgroup we may choose, in any of the many possible ways, a system of
representatives $S$ consisting of the numbers $\alpha_1, \alpha_2,
\ldots, \alpha_{\frac{p-1}{3}}$. What characterizes these
numbers is that the congruence $\alpha_i \equiv f^{\pm1}\alpha_k
\quad(\mathrm{mod.}\ p)$ cannot hold for any two indices $i, k$.
Now define the systems $S'$ and $S''$, each containing $\dfrac{p-1}{3}$ numbers
$\beta, \gamma$, by

\[
\beta_i \equiv f\alpha_i, \qquad \gamma_i \equiv f^2\alpha_i \quad(\mathrm{mod.}\ p),
\]
\noindent and the $(p-1)$ numbers $\alpha, \beta, \gamma$ exhaust precisely the complete
system of residues\hfill\\
$1, 2, \ldots, {p-1} \quad(\mathrm{mod.}\ p)$.

If $a$ is any number indivisible by $p$, then

\[
a^{\frac{p-1}{3}} \equiv f^{\lambda} \equiv \left[\frac{a}{p}\right] \quad(\mathrm{mod.}\ p)
\]
\noindent is the »cubic character« of $a$ with respect to $p$, where
$\lambda = 0, 1$ or $2$.

For a theorem on this cubic character, form the
$\dfrac{p-1}{3}$ products $a\alpha_1, a\alpha_2, \ldots
a\alpha_{\frac{p-1}{3}}$. These numbers evidently again form a system of
representatives, since the congruence $a\alpha_i \equiv
f^{\pm1}a\alpha_k$ would immediately imply $\alpha_i \equiv f^{\pm1}\alpha_k \quad(\mathrm{mod.}\ p)$.
If $\mu$ of the numbers $a\alpha_i$ belong to the system $S'$ and $\nu$
to the system $S''$, their product mod. $p$ is evidently congruent to
$f^{\mu+2\nu}. \alpha_1\alpha_2\ldots\alpha_{\frac{p-1}{3}}$. This gives the result

\[
\left[\frac{a}{p}\right] \equiv a^{\frac{p-1}{3}} \equiv f^{\mu+2\nu} \quad(\mathrm{mod.}\ p),
\]
\noindent which is the exact analogue of the well-known lemma on which \textsc{Gauss}
based his third proof of the law of quadratic reciprocity.

